% TOP_PROBLEM: {"id": 2302081, "title": "Research Problems in Function Theory — Problem 2.81", "category_id": 8, "category": {"id": 8, "name": "analysis", "display_name": "Analysis", "description": "Limits, continuity, calculus, and function theory.", "slug": "analysis", "order_index": 8, "created_at": "2026-07-31T15:26:25.671Z"}, "status": "open", "problem_number": "AMR-022-2081", "set_id": 13}
% FIRST_POSTED: 2026-10-02
% ENTRY_KIND: statement_only
% UnsolvedMath ID 2302081; source catalogue code AMR-022-2081
% !TeX program = lualatex
% Definition and sources only; this file is not an LLM solution attempt.
\documentclass[11pt,a4paper]{article}
\usepackage[margin=25mm]{geometry}
\usepackage{fontspec}
\setmainfont{lmroman10-regular.otf}[BoldFont=lmroman10-bold.otf,
 ItalicFont=lmroman10-italic.otf,BoldItalicFont=lmroman10-bolditalic.otf]
\usepackage{amsmath,amssymb,mathtools}
\usepackage{unicode-math}
\setmathfont{latinmodern-math.otf}
\AtBeginDocument{\let\setminus\smallsetminus}
\usepackage{xurl,hyperref}
\hypersetup{colorlinks=true,urlcolor=blue}
\setcounter{secnumdepth}{0}
\setlength{\parindent}{0pt}
\setlength{\parskip}{5pt}
\setlength{\emergencystretch}{3em}
\begin{document}
\section*{Research Problems in Function Theory — Problem 2.81}\label{2302081}
{\small UnsolvedMath ID 2302081 · AMR-022-2081 · Analysis · AMR Open Problem Lists}
\subsection{Definitions and mathematical statement}
For an integer $d\ge2$, let $\operatorname{Rat}_d$ be the space of degree-$d$ rational maps of the Riemann sphere. Represent a map by coprime homogeneous degree-$d$ polynomials, modulo common nonzero scaling; the induced coefficient topology is used. A critical point has local mapping degree greater than one. The Julia set $J(g)$ is the complement of the largest open set on which the iterates $\{g^n:n\ge0\}$ form a normal family in the spherical metric.
Assume $J(g)=\widehat{\mathbb C}$. Let $m$ denote spherical area measure, normalized if desired to have total mass one. Does every Borel set $B\subset\widehat{\mathbb C}$ satisfying
\[
g^{-1}(B)=B
\]
have either $m(B)=0$ or $m(\widehat{\mathbb C}\setminus B)=0$?
This is ergodicity of the map for the spherical-area measure class. The measure itself need not be preserved by $g$, so the question is not restricted to area-preserving maps or to invariance of an integral. Since a nonconstant rational map is nonsingular for area, invariant sets modulo null sets give the corresponding measure-class formulation.
The hypothesis that the Julia set fills the sphere removes the open Fatou components, but does not by itself rule out a measurable invariant partition with both pieces of positive area. The requested conclusion concerns every invariant Borel set. Ergodicity of a different canonical invariant probability measure, such as the measure of maximal entropy, is a separate statement and is not the question here.
\subsection{Short English statement}
Is a rational map with Julia set equal to the whole sphere ergodic for the spherical-area measure class?
\subsection{Sources}
\begin{itemize}
\item[S1] ULAM.AI and the UnsolvedMath Contributors, supplied problems.json, record 2302081 (AMR-022-2081), AMR Open Problem Lists. Catalogue identity and source transcription; CC BY 4.0.
\url{https://huggingface.co/datasets/ulamai/UnsolvedMath}.
\item[S2] Hayman - Research Problems in Function Theory (2018), Problem 2.81. Original source indexed by AMR-022-2081.
\url{https://arxiv.org/abs/1809.07200}.
\end{itemize}
\end{document}
